\documentclass[aspectratio=169,11pt]{beamer}

\usetheme{Madrid}
\usecolortheme{seahorse}
\setbeamertemplate{navigation symbols}{}

\usepackage[T1]{fontenc}
\usepackage[utf8]{inputenc}
\usepackage{amsmath}
\usepackage{booktabs}
\usepackage{xcolor}
\usepackage{listings}
\usepackage{fancyvrb}
\usepackage{minted}
\usepackage{hyperref}

\definecolor{CodeBg}{RGB}{248,248,248}
\definecolor{CodeFrame}{RGB}{220,220,220}
\definecolor{CodeKeyword}{RGB}{0,70,140}
\definecolor{CodeComment}{RGB}{92,112,92}
\definecolor{CodeString}{RGB}{150,60,30}
\definecolor{TraceDone}{RGB}{20,20,20}
\definecolor{TraceFuture}{RGB}{155,155,155}
\definecolor{TraceNote}{RGB}{82,82,82}
\definecolor{TerminalBg}{RGB}{36,40,44}
\definecolor{TerminalText}{RGB}{235,238,240}

\lstdefinestyle{python}{
  language=Python,
  basicstyle=\ttfamily\small,
  keywordstyle=\color{CodeKeyword}\bfseries,
  commentstyle=\color{CodeComment},
  stringstyle=\color{CodeString},
  showstringspaces=false,
  columns=fullflexible,
  keepspaces=true,
  frame=single,
  framerule=0.3pt,
  rulecolor=\color{CodeFrame},
  backgroundcolor=\color{CodeBg},
  breaklines=true,
  tabsize=2
}

\lstdefinestyle{mlir}{
  basicstyle=\ttfamily\scriptsize,
  columns=fullflexible,
  keepspaces=true,
  frame=single,
  framerule=0.3pt,
  rulecolor=\color{CodeFrame},
  backgroundcolor=\color{CodeBg},
  breaklines=true
}

\lstdefinestyle{traceir}{
  basicstyle=\ttfamily\tiny,
  columns=fullflexible,
  keepspaces=true,
  backgroundcolor=\color{CodeBg},
  breaklines=true,
  tabsize=2,
  escapeinside={(*@}{@*)}
}

\lstdefinestyle{terminal}{
  basicstyle=\ttfamily\tiny\color{TerminalText},
  columns=fullflexible,
  keepspaces=true,
  backgroundcolor=\color{TerminalBg},
  breaklines=true,
  tabsize=2
}

\DefineVerbatimEnvironment{TraceCode}{Verbatim}{
  fontsize=\scriptsize,
  commandchars=!\{\}
}

\setminted{
  fontsize=\small,
  breaklines,
  autogobble,
  frame=single,
  framesep=2mm
}

\newcommand{\RepoURL}{https://github.com/tongzhou8086/cute-tutorial/blob/main}
\newcommand{\code}[1]{\texttt{\detokenize{#1}}}
\newcommand{\codefile}[1]{\href{\RepoURL/#1}{\texttt{\detokenize{#1}}}}
\newcommand{\Cpp}{C\kern-.05em\texttt{+\kern-.02em+}}

\title[CuTe DSL Overview]{CuTe DSL Overview}
\author{cute-tutorial}
\date{\today}

\begin{document}

\begin{frame}[plain]
  \titlepage
\end{frame}

\begin{frame}{What is CuTe DSL?}
  \begin{center}
    \large
    CuTe DSL is a \textbf{Python} decorator-based DSL for writing\\
    hardware-aware \textbf{GPU kernels} with JIT compilation.
  \end{center}

  \vspace{1em}
  \begin{itemize}
    \item Evolved from \Cpp{} CUTLASS / CuTe
    \item Zero-cost abstractions with low-level hardware control
    \item DLPack interop, native types, and JIT caching
  \end{itemize}
\end{frame}

\begin{frame}{Outline --- Three Modules}
  \begin{itemize}
    \item[\textbf{1}] \hyperlink{mod:core}{\textbf{Core Foundations}}\\
          {\small Python $\rightarrow$ GPU code, layouts \& tensors}
    \vspace{0.6em}
    \item[\textbf{2}] \hyperlink{mod:simt}{\textbf{SIMT Programming Model}}\\
          {\small threads move data and do scalar work}
    \vspace{0.6em}
    \item[\textbf{3}] \hyperlink{mod:gemm}{\textbf{Hardware GEMM \& Async Pipelines}}\\
          {\small TMA, MMA, and mbarriers toward Blackwell GEMM}
  \end{itemize}
\end{frame}

\begin{frame}[plain]
  \hypertarget{mod:core}{}
  \centering
  \vfill
  {\small\color{gray} Module 1}\\[0.5em]
  {\Huge\bfseries Core Foundations}\\[1em]
  {\large Python $\rightarrow$ GPU code, layouts \& tensors}
  \vfill
\end{frame}

\begin{frame}{Core Foundations Roadmap}
  \begin{itemize}
    \item \textbf{Entry points:} \code{@cute.jit} for host-side JIT logic, \code{@cute.kernel} for GPU kernels.
    \vspace{0.45em}
    \item \textbf{Tracing model:} the same Python code builds IR first, then compiled code runs later.
    \vspace{0.45em}
    \item \textbf{Dynamic vs. constant arguments:} runtime values stay as IR arguments; \code{Constexpr} values specialize the generated code.
    \vspace{0.45em}
    \item \textbf{JIT cache and executors:} direct calls trace every time; \code{cute.compile} returns a reusable executor.
    \vspace{0.45em}
    \item \textbf{Tensor layouts:} CuTe describes tensors with \code{(Shape):(Stride)} and supports static or dynamic layout fields.
  \end{itemize}
\end{frame}

\begin{frame}{Two Decorators: \code{@cute.jit} and \code{@cute.kernel}}
Two entry points:
\begin{itemize}
  \item \code{@cute.jit}: host-side JIT functions
  \item \code{@cute.kernel}: GPU kernels
\end{itemize}

\end{frame}


\begin{frame}[fragile]{The First Example}
\begin{minted}{python}
import cutlass
import cutlass.cute as cute

@cute.jit
def add_dynamicexpr(b: cutlass.Float32):
    a: cutlass.Float32 = 2.0
    result = a + b
    print("[meta-stage] result =", result)  
    cute.printf("[object-stage] result = %f\n", result)  

add_dynamicexpr(5.0)

\end{minted}  
\end{frame}

\begin{frame}{General Mechanism: Same Python Code Runs Twice}
  \begin{center}
    \large
    The same function body is used in two different contexts:\\
    first to \textbf{build IR}, then to \textbf{run the compiled program}.
  \end{center}

  \vspace{0.6em}
  \begin{enumerate}
    \item \textbf{Pre-stage:} CuTe DSL rewrites the Python AST to capture structure.
    \item \textbf{Meta-stage:} Python runs on the host with proxy arguments; tensor operations emit MLIR.
    \item \textbf{Object-stage:} MLIR is lowered and compiled; emitted operations run at runtime.
  \end{enumerate}

  \vspace{0.4em}
  \begin{itemize}
    \item \code{print()} runs during tracing.
    \item \code{cute.printf()} does not print during tracing; it is emitted as \code{cute.print}.
  \end{itemize}

\end{frame}



\begin{frame}{What Happens When We Call \code{add_dynamicexpr(5.0)}?}
  \begin{center}
    \code{add_dynamicexpr(5.0)}
    \quad$\rightarrow$\quad
    \textbf{rewrite}
    \quad$\rightarrow$\quad
    \textbf{trace}
    \quad$\rightarrow$\quad
    \textbf{lower + compile}
  \end{center}

  \vspace{0.8em}
  \begin{enumerate}
    \item \textbf{Pre-stage: rewrite the Python AST}\\
          {\small Capture control-flow structure. For this straight-line example, there is no branch or loop to preserve.}
    \vspace{0.45em}
    \item \textbf{Meta-stage: run Python with proxy values (tracing)}\\
          {\small The function body executes on the host; tensor operations append MLIR to the current block.}
    \vspace{0.45em}
    \item \textbf{Object-stage: lower and compile MLIR}\\
          {\small The completed IR is lowered toward PTX/SASS; emitted operations run later.}
  \end{enumerate}

  \vfill
  \begin{center}
    \small The next slides zoom into the \textbf{meta-stage} one statement at a time.
  \end{center}
\end{frame}


\begin{frame}[fragile]{Tracing Step by Step: \code{add_dynamicexpr(5.0)}}
\begin{columns}[T,totalwidth=\textwidth]
  \begin{column}{0.48\textwidth}
    \small\textbf{Tracing code}
\begin{TraceCode}
!color{TraceDone}@cute.jit
!color{TraceDone}def add_dynamicexpr(b: cutlass.Float32):
!color{TraceDone}    a: cutlass.Float32 = 2.0
!color{TraceFuture}    result = a + b
!color{TraceFuture}    print("[meta-stage] result =", result)
!color{TraceFuture}    cute.printf("[object-stage] result = %f\n",
!color{TraceFuture}                result)
\end{TraceCode}
  \end{column}
  \begin{column}{0.52\textwidth}
    \small\textbf{MLIR}
\begin{lstlisting}[style=traceir]
// no IR operation yet
\end{lstlisting}

    \vspace{0.45em}
    \small\textbf{Terminal}
\begin{lstlisting}[style=terminal]
<no output yet>
\end{lstlisting}
  \end{column}
\end{columns}
\end{frame}

\begin{frame}[fragile]{Tracing Step by Step: \code{add_dynamicexpr(5.0)}}
\begin{columns}[T,totalwidth=\textwidth]
  \begin{column}{0.48\textwidth}
    \small\textbf{Tracing code}
\begin{TraceCode}
!color{TraceDone}@cute.jit
!color{TraceDone}def add_dynamicexpr(b: cutlass.Float32):
!color{TraceDone}    a: cutlass.Float32 = 2.0
!color{TraceDone}    result = a + b
!color{TraceFuture}    print("[meta-stage] result =", result)
!color{TraceFuture}    cute.printf("[object-stage] result = %f\n",
!color{TraceFuture}                result)
\end{TraceCode}
  \end{column}
  \begin{column}{0.52\textwidth}
    \small\textbf{MLIR}
\begin{lstlisting}[style=traceir]
%cst = arith.constant 2.000000e+00 : f32
%0 = arith.addf %cst, %arg0 : f32
\end{lstlisting}

    \vspace{0.45em}
    \small\textbf{Terminal}
\begin{lstlisting}[style=terminal]
<no output yet>
\end{lstlisting}
  \end{column}
\end{columns}
\end{frame}

\begin{frame}[fragile]{Tracing Step by Step: \code{add_dynamicexpr(5.0)}}
\begin{columns}[T,totalwidth=\textwidth]
  \begin{column}{0.48\textwidth}
    \small\textbf{Tracing code}
\begin{TraceCode}
!color{TraceDone}@cute.jit
!color{TraceDone}def add_dynamicexpr(b: cutlass.Float32):
!color{TraceDone}    a: cutlass.Float32 = 2.0
!color{TraceDone}    result = a + b
!color{TraceDone}    print("[meta-stage] result =", result)
!color{TraceFuture}    cute.printf("[object-stage] result = %f\n",
!color{TraceFuture}                result)
\end{TraceCode}
  \end{column}
  \begin{column}{0.52\textwidth}
    \small\textbf{MLIR}
\begin{lstlisting}[style=traceir]
%cst = arith.constant 2.000000e+00 : f32
%0 = arith.addf %cst, %arg0 : f32
\end{lstlisting}

    \vspace{0.45em}
    \small\textbf{Terminal}
\begin{lstlisting}[style=terminal]
[meta-stage] result = ?
\end{lstlisting}
  \end{column}
\end{columns}
\end{frame}

\begin{frame}[fragile]{Tracing Step by Step: \code{add_dynamicexpr(5.0)}}
\begin{columns}[T,totalwidth=\textwidth]
  \begin{column}{0.48\textwidth}
    \small\textbf{Tracing code}
\begin{TraceCode}
!color{TraceDone}@cute.jit
!color{TraceDone}def add_dynamicexpr(b: cutlass.Float32):
!color{TraceDone}    a: cutlass.Float32 = 2.0
!color{TraceDone}    result = a + b
!color{TraceDone}    print("[meta-stage] result =", result)
!color{TraceDone}    cute.printf("[object-stage] result = %f\n",
!color{TraceDone}                result)
\end{TraceCode}
  \end{column}
  \begin{column}{0.52\textwidth}
    \small\textbf{MLIR}
\begin{lstlisting}[style=traceir]
%cst = arith.constant 2.000000e+00 : f32
%0 = arith.addf %cst, %arg0 : f32
cute.print("[object-stage] result = %f\0A\0A", %0) : f32
\end{lstlisting}

    \vspace{0.45em}
    \small\textbf{Terminal}
\begin{lstlisting}[style=terminal]
[meta-stage] result = ?
\end{lstlisting}

    \vspace{0.25em}
    \footnotesize\color{TraceNote}
    No new trace-time output: \code{cute.printf} was emitted as \code{cute.print}.
  \end{column}
\end{columns}
\end{frame}

\begin{frame}[fragile]{Tracing Step by Step: \code{add_constexpr(5.0)}}
\begin{columns}[T,totalwidth=\textwidth]
  \begin{column}{0.48\textwidth}
    \small\textbf{Tracing code}
\begin{TraceCode}
!color{TraceDone}@cute.jit
!color{TraceDone}def add_constexpr(b: cutlass.Constexpr):
!color{TraceDone}    a: cutlass.Float32 = 2.0
!color{TraceDone}    result = a + b
!color{TraceDone}    print("[meta-stage] result =", result)
!color{TraceFuture}    cute.printf("[object-stage] result = %f\n",
!color{TraceFuture}                result)
\end{TraceCode}
  \end{column}
  \begin{column}{0.52\textwidth}
    \small\textbf{MLIR}
\begin{lstlisting}[style=traceir]
// no IR operation yet
\end{lstlisting}

    \vspace{0.45em}
    \small\textbf{Terminal}
\begin{lstlisting}[style=terminal]
[meta-stage] result = 7.0
\end{lstlisting}

    \vspace{0.25em}
    \footnotesize\color{TraceNote}
    Both inputs are known at trace time; \code{a} is explicitly typed as \code{Float32}.
  \end{column}
\end{columns}
\end{frame}

\begin{frame}[fragile]{Tracing Step by Step: \code{add_constexpr(5.0)}}
\begin{columns}[T,totalwidth=\textwidth]
  \begin{column}{0.48\textwidth}
    \small\textbf{Tracing code}
\begin{TraceCode}
!color{TraceDone}@cute.jit
!color{TraceDone}def add_constexpr(b: cutlass.Constexpr):
!color{TraceDone}    a: cutlass.Float32 = 2.0
!color{TraceDone}    result = a + b
!color{TraceDone}    print("[meta-stage] result =", result)
!color{TraceDone}    cute.printf("[object-stage] result = %f\n",
!color{TraceDone}                result)
\end{TraceCode}
  \end{column}
  \begin{column}{0.52\textwidth}
    \small\textbf{MLIR}
\begin{lstlisting}[style=traceir]
%cst = arith.constant 7.000000e+00 : f32
cute.print("[object-stage] result = %f\0A\0A", %cst) : f32
\end{lstlisting}

    \vspace{0.45em}
    \small\textbf{Terminal}
\begin{lstlisting}[style=terminal]
[meta-stage] result = 7.0
\end{lstlisting}

    \vspace{0.25em}
    \footnotesize\color{TraceNote}
    The generated IR has no runtime \texttt{\%arg0}; the value was baked in.
  \end{column}
\end{columns}
\end{frame}

\begin{frame}[fragile]{JIT Cache: Dynamic vs. \code{Constexpr} Values}
\begin{columns}[T,totalwidth=\textwidth]
  \begin{column}{0.49\textwidth}
    \small\textbf{Dynamic parameter}
\begin{TraceCode}
add_dynamicexpr(5.0)
add_dynamicexpr(8.0)
\end{TraceCode}

    \vspace{0.25em}
    \small\textbf{Generated IR shape}
\begin{lstlisting}[style=traceir]
%cst = arith.constant 2.000000e+00 : f32
%0 = arith.addf %cst, %arg0 : f32
cute.print(..., %0) : f32
\end{lstlisting}

    \vspace{0.2em}
    \footnotesize\color{TraceNote}
    The argument stays runtime data. New values with the same signature reuse
    the same native code because the IR shape is unchanged.
  \end{column}

  \begin{column}{0.49\textwidth}
    \small\textbf{\code{Constexpr} parameter}
\begin{TraceCode}
add_constexpr(5.0)  # result = 7.0
add_constexpr(8.0)  # result = 10.0
\end{TraceCode}

    \vspace{0.25em}
    \small\textbf{Generated IR changes}
\begin{lstlisting}[style=traceir]
// b = 5.0
%cst = arith.constant 7.000000e+00 : f32
cute.print(..., %cst) : f32

// b = 8.0
%cst = arith.constant 1.000000e+01 : f32
cute.print(..., %cst) : f32
\end{lstlisting}

    \vspace{0.2em}
    \footnotesize\color{TraceNote}
    The value is part of specialization. Changing it changes the IR hash,
    so the JIT cache must generate new native code.
  \end{column}
\end{columns}

\vspace{0.35em}
\begin{center}
  \footnotesize
  \textbf{Cache lookup uses the emitted MLIR hash, not just Python source.}\\
\end{center}
\end{frame}

\begin{frame}[fragile]{Control Flow: The Rewrite Becomes Visible}
\begin{columns}[T,totalwidth=\textwidth]
  \begin{column}{0.48\textwidth}
    \small\textbf{Source shape}
\begin{TraceCode}
@cute.jit
def choose_runtime(x: cutlass.Int32):
    if x > cutlass.Int32(0):
        print("[trace] building then branch")
        cute.printf("[runtime] positive\n")
    else:
        print("[trace] building else branch")
        cute.printf("[runtime] non-positive\n")
\end{TraceCode}
  \end{column}
  \begin{column}{0.52\textwidth}
    \small\textbf{Conceptual rewrite}
\begin{lstlisting}[style=traceir]
pred = x > cutlass.Int32(0)

if_generate(
    pred,
    then_branch,
    else_branch,
)
\end{lstlisting}

    \vspace{0.45em}
    \footnotesize\color{TraceNote}
    Straight-line code barely uses the pre-stage. Runtime control flow needs it:
    branch bodies become regions that can be traced into IR.
  \end{column}
\end{columns}
\end{frame}

\begin{frame}[fragile]{Runtime Branching: \code{if x > 0}}
\begin{columns}[T,totalwidth=\textwidth]
  \begin{column}{0.48\textwidth}
    \small\textbf{Tracing code}
\begin{TraceCode}
!color{TraceDone}@cute.jit
!color{TraceDone}def choose_runtime(x: cutlass.Int32):
!color{TraceDone}    if x > cutlass.Int32(0):
!color{TraceDone}        print("[trace] building then branch")
!color{TraceDone}        cute.printf("[runtime] positive\n")
!color{TraceDone}    else:
!color{TraceDone}        print("[trace] building else branch")
!color{TraceDone}        cute.printf("[runtime] non-positive\n")
\end{TraceCode}
  \end{column}
  \begin{column}{0.52\textwidth}
    \small\textbf{MLIR}
\begin{lstlisting}[style=traceir]
%c0_i32 = arith.constant 0 : i32
%0 = arith.cmpi sgt, %arg0, %c0_i32 : i32
scf.if %0 {
  cute.print("[runtime] positive\0A\0A")
} else {
  cute.print("[runtime] non-positive\0A\0A")
}
\end{lstlisting}

    \vspace{0.45em}
    \small\textbf{Terminal}
\begin{lstlisting}[style=terminal]
[trace] building then branch
[trace] building else branch
\end{lstlisting}

    \vspace{0.25em}
    \footnotesize\color{TraceNote}
    Both branch regions are traced; runtime chooses one branch.
  \end{column}
\end{columns}
\end{frame}

\begin{frame}[fragile]{Constexpr Value, Plain \code{if}: Still \code{scf.if}}
\begin{columns}[T,totalwidth=\textwidth]
  \begin{column}{0.48\textwidth}
    \small\textbf{Tracing code}
\begin{TraceCode}
!color{TraceDone}@cute.jit
!color{TraceDone}def choose_path(x: cutlass.Int32,
!color{TraceDone}                use_fast: cutlass.Constexpr):
!color{TraceDone}    if use_fast:
!color{TraceDone}        print("[trace] building fast branch")
!color{TraceDone}        cute.printf("[runtime] y = %d\n",
!color{TraceDone}                    x + cutlass.Int32(1))
!color{TraceDone}    else:
!color{TraceDone}        print("[trace] building slow branch")
!color{TraceDone}        cute.printf("[runtime] y = %d\n",
!color{TraceDone}                    x + cutlass.Int32(100))
\end{TraceCode}
  \end{column}
  \begin{column}{0.52\textwidth}
    \small\textbf{MLIR}
\begin{lstlisting}[style=traceir]
%true = arith.constant true
scf.if %true {
  %c1_i32 = arith.constant 1 : i32
  %0 = arith.addi %arg0, %c1_i32 : i32
  cute.print("[runtime] y = %d\0A\0A", %0) : i32
} else {
  %c100_i32 = arith.constant 100 : i32
  %0 = arith.addi %arg0, %c100_i32 : i32
  cute.print("[runtime] y = %d\0A\0A", %0) : i32
}
\end{lstlisting}

    \vspace{0.35em}
    \small\textbf{Terminal}
\begin{lstlisting}[style=terminal]
[trace] building fast branch
[trace] building slow branch
\end{lstlisting}

    \vspace{0.2em}
    \footnotesize\color{TraceNote}
    \code{Constexpr} alone does not make this a trace-time branch;
    plain \code{if} still uses the control-flow rewrite.
  \end{column}
\end{columns}
\end{frame}

\begin{frame}[fragile]{Trace-Time Branching: \code{if cutlass.const_expr(...)}}
\begin{columns}[T,totalwidth=\textwidth]
  \begin{column}{0.48\textwidth}
    \small\textbf{Tracing code}
\begin{TraceCode}
!color{TraceDone}@cute.jit
!color{TraceDone}def choose_path(x: cutlass.Int32,
!color{TraceDone}                use_fast: cutlass.Constexpr):
!color{TraceDone}    if cutlass.const_expr(use_fast):
!color{TraceDone}        print("[trace] building fast branch")
!color{TraceDone}        cute.printf("[runtime] y = %d\n",
!color{TraceDone}                    x + cutlass.Int32(1))
!color{TraceFuture}    else:
!color{TraceFuture}        print("[trace] building slow branch")
!color{TraceFuture}        cute.printf("[runtime] y = %d\n",
!color{TraceFuture}                    x + cutlass.Int32(100))
\end{TraceCode}
  \end{column}
  \begin{column}{0.52\textwidth}
    \small\textbf{MLIR}
\begin{lstlisting}[style=traceir]
%c1_i32 = arith.constant 1 : i32
%0 = arith.addi %arg0, %c1_i32 : i32
cute.print("[runtime] y = %d\0A\0A", %0) : i32
\end{lstlisting}

    \vspace{0.45em}
    \small\textbf{Terminal}
\begin{lstlisting}[style=terminal]
[trace] building fast branch
\end{lstlisting}

    \vspace{0.25em}
    \footnotesize\color{TraceNote}
    \code{cutlass.const_expr(...)} is the explicit signal to skip the control-flow rewrite.
    The predicate must be compile-time known; otherwise tracing errors.
  \end{column}
\end{columns}
\end{frame}

\begin{frame}[fragile]{Compile-Time Metaprogramming: Optional ReLU}
\begin{columns}[T,totalwidth=\textwidth]
  \begin{column}{0.52\textwidth}
    \small\textbf{Kernel sketch}
\begin{TraceCode}
@cute.kernel
def gemm(...,
         do_relu: cutlass.Constexpr):
    # main GEMM work
    ...

    if cutlass.const_expr(do_relu):
        # ReLU epilogue is emitted only
        # when do_relu is True
        ...
\end{TraceCode}
  \end{column}

  \begin{column}{0.46\textwidth}
    \small\textbf{Two specializations}
\begin{lstlisting}[style=traceir]
gemm(..., False)

// generated code:
main GEMM
store accumulator


gemm(..., True)

// generated code:
main GEMM
relu accumulator
store accumulator
\end{lstlisting}

    \vspace{0.25em}
    \footnotesize\color{TraceNote}
    The flag is a constant argument. It selects code during tracing, so the
    final kernel has no runtime branch for this option.
  \end{column}
\end{columns}
\end{frame}

\begin{frame}[fragile]{JIT Gotcha: Direct \code{@cute.jit} Calls}
\begin{columns}[T,totalwidth=\textwidth]
  \begin{column}{0.58\textwidth}
    \small\textbf{Code}
\begin{TraceCode}
import cutlass
import cutlass.cute as cute

@cute.jit
def hello():
    print("trace: hello")
    cute.printf("runtime: hello\n")

hello()
hello()
\end{TraceCode}
  \end{column}

  \begin{column}{0.38\textwidth}
    \small\textbf{Question}

    \vspace{0.6em}
    How many times does each line print?

    \vspace{0.7em}
    \begin{itemize}
      \item \code{trace: hello}
      \item \code{runtime: hello}
    \end{itemize}
  \end{column}
\end{columns}
\end{frame}


\begin{frame}[fragile]{Implicit Cache: Trace Every Call, Compile on Miss}
\begin{columns}[T,totalwidth=\textwidth]
  \begin{column}{0.48\textwidth}
    \small\textbf{First direct call}
\begin{lstlisting}[style=traceir]
hello()

trace Python
emit MLIR
cache lookup: miss
compile MLIR -> native code
run native code
\end{lstlisting}
  \end{column}

  \begin{column}{0.48\textwidth}
    \small\textbf{Second direct call}
\begin{lstlisting}[style=traceir]
hello()

trace Python
emit MLIR
cache lookup: hit
reuse compiled executor
run native code
\end{lstlisting}
  \end{column}
\end{columns}

\vspace{0.55em}
\begin{center}
  \footnotesize
  Because generated MLIR bytecode is part of the implicit cache key, CuTe must
  trace before it can decide whether compilation is a cache hit.
\end{center}
\end{frame}


\begin{frame}[fragile]{Use \code{cute.compile} to Reuse a JIT Executor}
\begin{columns}[T,totalwidth=\textwidth]
  \begin{column}{0.58\textwidth}
    \small\textbf{Code}
\begin{TraceCode}
import cutlass
import cutlass.cute as cute

@cute.jit
def hello():
    print("trace: hello")
    cute.printf("runtime: hello\n")

compiled_func = cute.compile(hello)

compiled_func()
compiled_func()
\end{TraceCode}
  \end{column}

  \begin{column}{0.38\textwidth}
    \small\textbf{Terminal}
\begin{lstlisting}[style=terminal]
trace: hello
runtime: hello
runtime: hello
\end{lstlisting}

    \vspace{0.35em}
    \footnotesize\color{TraceNote}
    In this pattern, \code{cute.compile} traces and compiles once, returning a
    JIT Executor. Reusing that executor skips Python tracing; calling
    \code{cute.compile} again would compile again.
  \end{column}
\end{columns}
\end{frame}


\begin{frame}[fragile]{Exercise: \code{Constexpr} and \code{cute.compile}}
\begin{columns}[T,totalwidth=\textwidth]
  \begin{column}{0.58\textwidth}
    \small\textbf{Code}
\begin{TraceCode}
import cutlass
import cutlass.cute as cute

@cute.jit
def add(a, b, print_result: cutlass.Constexpr):
    if print_result:
        cute.printf("Result: %d\n", a + b)
    return a + b

jit_executor = cute.compile(add, 1, 2, True)
jit_executor(1, 2)
\end{TraceCode}
  \end{column}

  \begin{column}{0.38\textwidth}
    \small\textbf{Questions}

    \vspace{0.6em}
    What prints during \code{cute.compile(...)}?

    \vspace{0.7em}
    What prints during \code{jit_executor(1, 2)}?

    \vspace{0.7em}
    Does the emitted MLIR keep control flow, or is it already specialized?

    \vspace{0.45em}
    \footnotesize\color{TraceNote}
    \code{print_result} is a compile-time argument, so it is not passed to
    the executor at runtime.
  \end{column}
\end{columns}
\end{frame}


\begin{frame}[fragile]{CuTe Tensor Layout: \code{(Shape):(Stride)}}
\begin{columns}[T,totalwidth=\textwidth]
  \begin{column}{0.46\textwidth}
    \small\textbf{Layout notation}
\begin{lstlisting}[style=traceir]
(3):(1)
(2,2):(2,1)
(?,?):(?{i64},1)
\end{lstlisting}
  \end{column}

  \begin{column}{0.54\textwidth}
    \small\textbf{Row-major example}
\begin{lstlisting}[style=traceir]
layout = (M,N):(N,1)

offset(i,j) = i * N + j
\end{lstlisting}
  \end{column}
\end{columns}

\vspace{1.3em}
\begin{center}
  \small
  A layout maps logical coordinates to offsets.\\[0.55em]
  \footnotesize
  \code{Shape}: logical extents
  \quad
  \code{Stride}: offset step per mode
  \quad
  \code{?}: runtime-known dimension or stride
\end{center}
\end{frame}

\begin{frame}[fragile]{Argument Annotations: What Gets Inferred?}
\begin{columns}[T,totalwidth=\textwidth]
  \begin{column}{0.48\textwidth}
    \small\textbf{Unannotated scalars}
\begin{lstlisting}[style=traceir]
@cute.jit
def foo(x):
    cute.printf("x = {}\n", x)

foo(1.0)   // x: Float32, MLIR f32
foo(1)     // x: Int32,   MLIR i32
foo(True)  // x: Boolean, MLIR i1
\end{lstlisting}

    \vspace{0.25em}
    \footnotesize\color{TraceNote}
    By default, arguments are dynamic and CuTe infers their DSL types from the
    call-site values.
  \end{column}

  \begin{column}{0.48\textwidth}
    \small\textbf{Annotate when intent matters}
\begin{lstlisting}[style=traceir]
def foo(x: cutlass.Float64):
    ...

def bar(n: cutlass.Int64):
    ...

def baz(flag: cutlass.Constexpr):
    ...

def qux(tensor: cute.Tensor):
    ...
\end{lstlisting}

    \vspace{0.25em}
    \footnotesize\color{TraceNote}
    Use annotations for precision or width, compile-time constants, and type
    safety. For tensors, static vs. dynamic layout is decided by the argument
    object passed at the call site.
  \end{column}
\end{columns}
\end{frame}

\begin{frame}[fragile]{Static Layout: Wrap with \code{from_dlpack}}
\begin{columns}[T,totalwidth=\textwidth]
  \begin{column}{0.49\textwidth}
    \small\textbf{Host code}
\begin{lstlisting}[style=traceir,basicstyle=\ttfamily\scriptsize]
import cutlass.cute as cute
from cutlass.cute.runtime import from_dlpack

@cute.jit
def foo(tensor):
    print(tensor.layout)
    cute.printf("tensor: {}", tensor)

a = torch.tensor([1, 2, 3],
                 dtype=torch.uint16)
a_pack = from_dlpack(a)

compiled = cute.compile(foo, a_pack)
compiled(a_pack)
\end{lstlisting}
  \end{column}

  \begin{column}{0.49\textwidth}
    \small\textbf{Trace-time output}
\begin{lstlisting}[style=terminal]
(3):(1)
\end{lstlisting}

    \vspace{0.4em}
    \small\textbf{Runtime output}
\begin{lstlisting}[style=terminal]
tensor: ... o (3):(1) =
  ( 1, 2, 3 )
\end{lstlisting}

    \vspace{0.25em}
    \footnotesize\color{TraceNote}
    \code{print(tensor.layout)} runs while tracing. \code{cute.printf(...)}
    is emitted and prints when the compiled executor runs.
  \end{column}
\end{columns}
\end{frame}

\begin{frame}[fragile]{Dynamic Layout: Pass \code{torch.Tensor} Directly}
\begin{columns}[T,totalwidth=\textwidth]
  \begin{column}{0.49\textwidth}
    \small\textbf{Host code}
\begin{TraceCode}
a = torch.tensor([[1, 2], [3, 4]],
                 dtype=torch.uint16)
compiled = cute.compile(foo, a)
compiled(a)

b = torch.tensor([[11, 12],
                  [13, 14],
                  [15, 16]],
                 dtype=torch.uint16)
compiled(b)
\end{TraceCode}
  \end{column}

  \begin{column}{0.49\textwidth}
    \small\textbf{Trace-time layout}
\begin{lstlisting}[style=terminal]
(?,?):(?{i64},1)
\end{lstlisting}

    \vspace{0.45em}
    \small\textbf{Runtime layouts}
\begin{lstlisting}[style=terminal]
tensor: ... o (2,2):(2,1)
tensor: ... o (3,2):(2,1)
\end{lstlisting}

    \vspace{0.25em}
    \footnotesize\color{TraceNote}
    The generated code keeps selected layout fields dynamic, so one compiled
    function can handle multiple compatible runtime shapes.
  \end{column}
\end{columns}
\end{frame}

\begin{frame}[plain]
  \hypertarget{mod:simt}{}
  \centering
  \vfill
  {\small\color{gray} Module 2}\\[0.5em]
  {\Huge\bfseries SIMT Programming Model}\\[1em]
  {\large threads, CTAs, scalar work, and tile ownership}
  \vfill
\end{frame}

\begin{frame}{SIMT Module Roadmap}
  \begin{itemize}
    \item \textbf{Vector add:} first \code{@cute.kernel}; one thread owns one element.
    \vspace{0.45em}
    \item \textbf{Reductions:} threads cooperate on scalar work.
    \vspace{0.45em}
    \item \textbf{2D indexing:} map a flat thread id to \code{(row, col)} tensor coordinates.
    \vspace{0.45em}
    \item \textbf{CTA tile views:} use \code{cute.local_tile} so each CTA owns a rectangular tile.
    \vspace{0.45em}
    \item \textbf{Shared-memory staging:} introduce \code{make_layout}, \code{SmemAllocator}, and CTA barriers.
  \end{itemize}

\end{frame}

\begin{frame}[fragile]{Vector Add: Familiar CUDA-Style SIMT}
\begin{columns}[T,totalwidth=\textwidth]
  \begin{column}{0.52\textwidth}
    \small\textbf{GPU kernel}
\begin{TraceCode}
@cute.kernel
def vector_add_kernel(a: cute.Tensor,
                      b: cute.Tensor,
                      c: cute.Tensor):
    tidx, _, _ = cute.arch.thread_idx()
    bidx, _, _ = cute.arch.block_idx()
    bdim, _, _ = cute.arch.block_dim()

    idx = bidx * bdim + tidx
    n = a.shape[0]

    if idx < n:
        c[idx] = a[idx] + b[idx]
\end{TraceCode}
  \end{column}

  \begin{column}{0.46\textwidth}
    \small\textbf{Host JIT launcher}
\begin{TraceCode}
@cute.jit
def vector_add(a: cute.Tensor,
               b: cute.Tensor,
               c: cute.Tensor):
    n = a.shape[0]
    grid = cute.ceil_div(n, THREADS_PER_BLOCK)
    vector_add_kernel(a, b, c).launch(
        grid=(grid, 1, 1),
        block=(THREADS_PER_BLOCK, 1, 1),
    )
\end{TraceCode}

    \vspace{0.45em}
    \footnotesize\color{TraceNote}
    One thread owns one element. \code{cute.arch.*} gives CUDA-style indices;
    tensor indexing expresses global-memory loads and stores.
  \end{column}
\end{columns}
\end{frame}

\begin{frame}[fragile]{Reduction: Shared Memory + One Atomic}
\begin{columns}[T,totalwidth=\textwidth]
  \begin{column}{0.52\textwidth}
    \small\textbf{GPU kernel}
\begin{lstlisting}[style=traceir]
@cute.kernel
def reduce_kernel(src, out, smem_layout):
    tidx, _, _ = cute.arch.thread_idx()
    bidx, _, _ = cute.arch.block_idx()
    idx = bidx * BLOCK + tidx

    buf = SmemAllocator().allocate_tensor(
        cutlass.Float32, smem_layout)

    buf[tidx] = cutlass.Float32(0.0)
    if idx < src.shape[0]:
        buf[tidx] = src[idx]
    cute.arch.barrier()

    stride = BLOCK // 2
    while stride > 0:
        if tidx < stride:
            buf[tidx] += buf[tidx + stride]
        cute.arch.barrier()
        stride = stride // 2

    if tidx == 0:
        cute.arch.atomic_add(out.iterator, buf[0])
\end{lstlisting}
  \end{column}

  \begin{column}{0.46\textwidth}
    \small\textbf{Host JIT launcher}
\begin{TraceCode}
@cute.jit
def reduce_sum(src: cute.Tensor,
               out: cute.Tensor):
    n = src.shape[0]
    grid = cute.ceil_div(n, BLOCK)
    smem_layout = cute.make_layout(BLOCK)

    reduce_kernel(src, out, smem_layout).launch(
        grid=(grid, 1, 1),
        block=(BLOCK, 1, 1),
    )
\end{TraceCode}

    \vspace{0.45em}
    \footnotesize\color{TraceNote}
    The host initializes \code{out[0] = 0}. Each CTA reduces into shared memory;
    one thread atomically adds the CTA sum to the global result.
  \end{column}
\end{columns}
\end{frame}

\begin{frame}[fragile]{\code{cute.make_layout}: Shape + Stride}
\begin{columns}[T,totalwidth=\textwidth]
  \begin{column}{0.48\textwidth}
    \small\textbf{Signature}
\begin{lstlisting}[style=traceir]
cute.make_layout(shape, *, stride=None)
\end{lstlisting}

    \vspace{0.35em}
    \begin{itemize}
      \item \code{shape}: integer or tuple of mode sizes.
      \item \code{stride}: optional keyword-only stride.
      \item returns a \code{cute.Layout}.
    \end{itemize}
  \end{column}

  \begin{column}{0.48\textwidth}
    \small\textbf{Common uses}
\begin{lstlisting}[style=traceir]
// 1D compact scratch
smem_layout = cute.make_layout(BLOCK)

// 2D row-major tile
tile_layout = cute.make_layout(
    (TILE_M, TILE_N),
    stride=(TILE_N, 1),
)

smem = SmemAllocator()
buf = smem.allocate_tensor(
    cutlass.Float32, smem_layout)
\end{lstlisting}

    \vspace{0.25em}
    \footnotesize\color{TraceNote}
    If \code{stride} is omitted, CuTe computes a default "left-most" layout.
    Pass \code{stride=} when the physical memory order matters.
  \end{column}
\end{columns}

\vfill
{\tiny API doc: \url{https://docs.nvidia.com/cutlass/latest/media/docs/pythonDSL/cute_dsl_api/cute.html\#cutlass.cute.make_layout}}

\end{frame}

\begin{frame}[fragile]{2D Indexing: Flat Thread Id to \code{(row, col)}}
\begin{columns}[T,totalwidth=\textwidth]
  \begin{column}{0.52\textwidth}
    \small\textbf{GPU kernel}
\begin{lstlisting}[style=traceir]
@cute.kernel
def copy_2d_kernel(src: cute.Tensor,
                   dst: cute.Tensor):
    tidx, _, _ = cute.arch.thread_idx()
    bidx, _, _ = cute.arch.block_idx()
    bdim, _, _ = cute.arch.block_dim()

    linear = bidx * bdim + tidx
    rows, cols = src.shape
    numel = rows * cols

    if linear < numel:
        row = linear // cols
        col = linear % cols
        dst[row, col] = src[row, col]
\end{lstlisting}
  \end{column}

  \begin{column}{0.46\textwidth}
    \small\textbf{Host JIT launcher}
\begin{lstlisting}[style=traceir]
@cute.jit
def copy_2d(src: cute.Tensor,
            dst: cute.Tensor):
    rows, cols = src.shape
    n = rows * cols
    grid = cute.ceil_div(
        n, THREADS_PER_BLOCK)

    copy_2d_kernel(src, dst).launch(
        grid=(grid, 1, 1),
        block=(THREADS_PER_BLOCK, 1, 1),
    )
\end{lstlisting}

    \vspace{0.45em}
    \footnotesize\color{TraceNote}
    Same SIMT skeleton as vector add. The only new step is decoding one flat
    element id into tensor coordinates.
  \end{column}
\end{columns}
\end{frame}

\begin{frame}[fragile]{CTA Tile Views with \code{cute.local_tile}}
\begin{columns}[T,totalwidth=\textwidth]
  \begin{column}{0.56\textwidth}
    \small\textbf{GPU kernel}
\begin{lstlisting}[style=traceir]
@cute.kernel
def local_tile_copy_kernel(src: cute.Tensor,
                           dst: cute.Tensor):
    tidx, _, _ = cute.arch.thread_idx()
    block_m, block_n, _ = cute.arch.block_idx()
    block_dim, _, _ = cute.arch.block_dim()

    src_tile = cute.local_tile(
        src, (TILE_M, TILE_N), (block_m, block_n))
    dst_tile = cute.local_tile(
        dst, (TILE_M, TILE_N), (block_m, block_n))

    i = tidx
    while i < TILE_M * TILE_N:
        row = i // TILE_N
        col = i % TILE_N
        dst_tile[row, col] = src_tile[row, col]
        i += block_dim
\end{lstlisting}
  \end{column}

  \begin{column}{0.40\textwidth}
    \small\textbf{Host JIT launcher}
\begin{lstlisting}[style=traceir]
@cute.jit
def local_tile_copy(src: cute.Tensor,
                    dst: cute.Tensor):
    rows, cols = src.shape
    grid_m = cute.ceil_div(rows, TILE_M)
    grid_n = cute.ceil_div(cols, TILE_N)

    local_tile_copy_kernel(src, dst).launch(
        grid=(grid_m, grid_n, 1),
        block=(THREADS_PER_CTA, 1, 1),
    )
\end{lstlisting}

    \vspace{0.35em}
    \footnotesize\color{TraceNote}
    \textbf{\code{local_tile} returns a tensor view; it does not copy data.}
  \end{column}
\end{columns}
\end{frame}

\end{document}
